\section{螺旋度与横率分布}
在前几节中，我们考虑了非极化夸克分布的匹配条件。本节给出夸克螺旋度分布与横率分布的结果。对于纵向极化夸克中的夸克螺旋度分布，把式 (1) 中的 $\gamma^z$ 替换为 $\gamma^z\gamma^5$ 即可得到准分布 $\Delta \tilde q^{(1)}(x)$。单圈结果为
\begin{align}\label{helivertexnoexp}
\Delta \tilde q^{(1)}(x)&=\frac{\alpha_S C_F}{2\pi} \left\{ \begin{array} {ll} \frac{1+x^2}{1-x}\ln \frac{x(\Lambda(x)-x P^z)}{(x-1)(\Lambda(1-x)+P^z(1-x))}+1-\frac{x P^z}{\Lambda(x)}+\frac{x\Lambda(1-x)+(1-x)\Lambda(x)}{(1-x)^2 P^z}\ , & x>1\ , \\
\ \\
\frac{1+x^2}{1-x}\ln\frac{(P^z)^2}{m^2}+\frac{1+x^2}{1-x}\ln \frac{4x(\Lambda(x)-x P^z)}{(1-x)(\Lambda(1-x)+(1-x)P^z)}-\frac{4}{1-x}+2x+3\ & \\
-\frac{x P^z}{\Lambda(x)}+\frac{x\Lambda(1-x)+(1-x)\Lambda(x)}{(1-x)^2 P^z}\ , & 0<x<1\ . \\
\ \\
\frac{1+x^2}{1-x}\ln \frac{(x-1)(\Lambda(x)-x P^z)}{x(\Lambda(1-x)+(1-x)P^z)}-1-\frac{x P^z}{\Lambda(x)}+\frac{x\Lambda(1-x)+(1-x)\Lambda(x)}{(1-x)^2P^z}\ , & x<0\ . \end{array} \right.
\end{align}
取极限 $\mu\rightarrow \infty$ 得到
\begin{align}\label{helivertexLgmu}
\Delta \tilde q^{(1)}(x)&=\frac{\alpha_S C_F}{2\pi} \left\{ \begin{array} {ll} \frac{1+x^2}{1-x}\ln \frac{x}{x-1}+1+\frac{\mu}{(1-x)^2P^z}\ , & x>1\ , \\ \frac{1+x^2}{1-x}\ln \frac{(P^z)^2}{m^2}+\frac{1+x^2}{1-x}\ln\frac{4x}{1-x}-\frac{4}{1-x}+2x+3+\frac{\mu}{(1-x)^2P^z}\ , & 0<x<1\ , \\ \frac{1+x^2}{1-x}\ln \frac{x-1}{x}-1+\frac{\mu}{(1-x)^2P^z}\ , & x<0\ . \end{array} \right.
\end{align}
光锥分布的结果同样由取 $P^z\to\infty$ 得到
\begin{align}\label{helivertexIMF}
\Delta \tilde q^{(1)}(x)&=\frac{\alpha_S C_F}{2\pi} \left\{ \begin{array} {ll} 0\ , & x>1\  \mbox{或}\  x<0\ , \\ \frac{1+x^2}{1-x}\ln \frac{\mu^2}{m^2}-\frac{1+x^2}{1-x}\ln{(1-x)^2}-\frac{2}{1-x}+2x\ , & 0<x<1\ .\end{array} \right.
\end{align}
注意，与非极化情形一样，准夸克螺旋度分布中的共线奇异性与光锥分布中的完全相同。


类似地，对于横向极化夸克中的横率分布，把式 (1) 中的 $\gamma^z$ 替换为 $\gamma^z\gamma^\perp\gamma^5$ 即可得到准分布 $\delta \tilde q^{(1)}(x)$。单圈结果为
\begin{align}\label{transvertexnoexp}
\delta \tilde q^{(1)}(x)&=\frac{\alpha_S C_F}{2\pi} \left\{ \begin{array} {ll} \frac{2x}{1-x}\ln \frac{x(\Lambda(x)-xP^z)}{(x-1)(\Lambda(1-x)+P^z(1-x))}+\frac{x\Lambda(1-x)+(1-x)\Lambda(x)}{(1-x)^2 P^z}\ , & x>1\ , \\
\ \\
\frac{2x}{1-x}\ln\frac{(P^z)^2}{m^2}+\frac{2x}{1-x}\ln \frac{4x(\Lambda(x)-x P^z)}{(1-x)(\Lambda(1-x)+(1-x)P^z)}-\frac{4x}{1-x}+\frac{x\Lambda(1-x)+(1-x)\Lambda(x)}{(1-x)^2 P^z}\ , & 0<x<1\ , \\
\ \\
\frac{2x}{1-x}\ln \frac{(x-1)(\Lambda(x)-x P^z)}{x(\Lambda(1-x)+(1-x)P^z)}+\frac{x\Lambda(1-x)+(1-x)\Lambda(x)}{(1-x)^2P^z}\ , & x<0\ . \end{array} \right.
\end{align}
取极限 $\mu\rightarrow \infty$ 给出
\begin{align}\label{transvertexLgmu}
\delta \tilde q^{(1)}(x)&=\frac{\alpha_S C_F}{2\pi} \left\{ \begin{array} {ll} \frac{2x}{1-x}\ln \frac{x}{x-1}+\frac{\mu}{(1-x)^2P^z}\ , & x>1\ , \\
\frac{2x}{1-x}\ln \frac{(P^z)^2}{m^2}+\frac{2x}{1-x}\ln\frac{4x}{1-x}-\frac{4x}{1-x}+\frac{\mu}{(1-x)^2P^z}\ , & 0<x<1\ ,\\
\frac{2x}{1-x}\ln \frac{x-1}{x}+\frac{\mu}{(1-x)^2P^z}\ , & x<0\ , \end{array} \right.
\end{align}
而无限动量系中的结果则为
\begin{align}\label{transvertexIMF}
\delta q^{(1)}(x)&=\frac{\alpha_S C_F}{2\pi} \left\{ \begin{array} {ll} 0\ , & x>1\  \mbox{或}\  x<0\ , \\ \frac{2x}{1-x}\ln \frac{\mu^2}{m^2}-\frac{2x}{1-x}\ln{(1-x)^2}-\frac{2x}{1-x}\ , & 0<x<1\ .\end{array} \right.
\end{align}

可以像在式 (12) 中那样，为螺旋度分布和横率分布构造类似的匹配条件。
注意到夸克自能与之前相同，我们在此只列出匹配因子的结果。
对夸克螺旋度分布，$\xi>1$ 时有
\begin{equation}
   \Delta Z^{(1)}(\xi)/C_F = \left(\frac{1+\xi^2}{1-\xi}\right)\ln \frac{\xi}{\xi-1} + 1 + \frac{1}{(1-\xi)^2}\frac{\mu}{P^z} \ ,
\end{equation}
而对 $0<\xi<1$
\begin{equation}
   \Delta Z^{(1)}(\xi)/C_F = \left(\frac{1+\xi^2}{1-\xi}\right)\ln \frac{(P^z)^2}{\mu^2}+\left(\frac{1+\xi^2}{1-\xi}\right)\ln\big[{4\xi(1-\xi)}\big]-\frac{2}{1-\xi}+3 +\frac{\mu}{(1-\xi)^2P^z} \ ,
\end{equation}
对 $\xi<0$
\begin{equation}
   \Delta Z^{(1)}(\xi)/C_F = \left(\frac{1+\xi^2}{1-\xi}\right)\ln \frac{\xi-1}{\xi}-1 +\frac{\mu}{(1-\xi)^2P^z} \ .
\end{equation}
线性发散项与非极化情形中的相同。

最后，在横率分布的因子化定理中，$\xi>1$ 时的匹配因子为
\begin{equation}
   \delta Z^{(1)}(\xi)/C_F = \left(\frac{2\xi}{1-\xi}\right)\ln \frac{\xi}{\xi-1} + \frac{1}{(1-\xi)^2}\frac{\mu}{P^z} \ ,
\end{equation}
而对 $0<\xi<1$
\begin{equation}
   \delta Z^{(1)}(\xi)/C_F = \left(\frac{2\xi}{1-\xi}\right)\ln \frac{(P^z)^2}{\mu^2}+\left(\frac{2\xi}{1-\xi}\right)\ln\big[{4\xi(1-\xi)}\big]-\frac{2\xi}{1-\xi} +\frac{\mu}{(1-\xi)^2P^z} \ ,
\end{equation}
对 $\xi<0$
\begin{equation}
   \delta Z^{(1)}(\xi)/C_F = \left(\frac{2\xi}{1-\xi}\right)\ln \frac{\xi-1}{\xi}+\frac{\mu}{(1-\xi)^2P^z} \ .
\end{equation}
这里同样存在一个线性发散的贡献。在 $\xi=1$ 附近，与前一样需要计入来自自能的额外贡献。
